% !TEX root = ../main.tex
\chapter{Étude préalable}

% ─────────────────────────────────────────────────────────────
\section{Introduction}

Ce chapitre constitue le socle du projet ASM Track. Nous y présentons l'organisme d'accueil,
définissons le cadre général du projet, analysons les solutions existantes sur le marché,
exposons la problématique identifiée, décrivons la solution proposée, puis détaillons la
méthodologie de travail adoptée tout au long du développement.

% ─────────────────────────────────────────────────────────────
\section{Présentation de l'organisme d'accueil}

Étant une Entreprise de Services du Numérique (ESN), \textbf{All Soft Multimédia (ASM)}
est une structure souple et réactive. Elle propose des solutions sur mesure dans le domaine
des nouvelles technologies et des systèmes d'information.

\begin{figure}[H]
  \centering
  \includegraphics[width=4cm]{asm-logo}
  \caption{Logo ASM}
  \label{fig:asm-logo}
\end{figure}

Créée en \textbf{2007}, ASM est une société spécialisée dans l'ingénierie logicielle, la
création de sites web, le développement d'applications mobiles ainsi que les systèmes de
points de vente. Elle est présente en Tunisie à \textbf{Sfax}, \textbf{Sousse} et \textbf{Tunis},
en Algérie à \textbf{Alger}, et au Sénégal à \textbf{Dakar}.

% ─────────────────────────────────────────────────────────────
\section{Cadre général du projet}

Ce projet de fin d'études, réalisé au sein d'ASM, s'inscrit dans le cadre de l'obtention
du \textbf{Diplôme National d'Ingénieur en Informatique} à l'Institut Supérieur d'Informatique
et de Multimédia de Sfax (ISIMS).

L'objectif est de concevoir et de développer \textbf{ASM Track}, une plateforme SaaS
multi-tenant de gestion et de suivi des livraisons en temps réel. Le projet couvre
l'intégralité de la chaîne logistique du dernier kilomètre, depuis l'importation des commandes
depuis un ERP jusqu'à la livraison finale avec capture de preuve, en passant par la
planification des tournées et le suivi en temps réel.

% ─────────────────────────────────────────────────────────────
\section{Contexte et problématique}

Pour illustrer notre propos, nous menons dans ce qui suit une étude contextuelle partant
de l'analyse des solutions de gestion des livraisons présentes sur le marché, allant vers
leur critique, pour aboutir à la description de notre problématique.

\subsection{Études des applications existantes sur le marché}

Avant d'entamer notre projet, il est essentiel de réaliser une veille externe. C'est pourquoi
il est important d'examiner les applications déjà présentes sur le marché afin de mieux
comprendre les aspects relatifs à notre contexte et d'identifier les manques auxquels notre
solution doit répondre.

\subsubsection*{Bringgg}

\begin{wrapfigure}{r}{3.5cm}
  \centering
  \includegraphics[width=3cm]{bringgg-logo}
  \caption{Logo Bringgg}
\end{wrapfigure}

Bringgg est une plateforme SaaS de gestion des livraisons du dernier kilomètre, fondée en
Israël et opérant à l'échelle internationale. Elle propose une suite complète d'outils pour
les opérateurs logistiques : planification des tournées, affectation des chauffeurs, suivi
GPS en temps réel et notifications automatiques aux destinataires. La plateforme se distingue
par son interface intuitive et ses capacités d'intégration avec des systèmes de e-commerce
tels que Shopify et WooCommerce, ce qui lui permet de répondre aux besoins des entreprises
de distribution à grande échelle.

\subsubsection*{Onfleet}

\begin{wrapfigure}{r}{3.5cm}
  \centering
  \includegraphics[width=3cm]{onfleet-logo}
  \caption{Logo Onfleet}
\end{wrapfigure}

Onfleet est une solution américaine de gestion des livraisons destinée aux entreprises de
taille moyenne et grande. Elle offre des fonctionnalités avancées de dispatch automatique,
de suivi en temps réel, d'analyse des performances et de communication entre les
dispatchers et les chauffeurs via une application mobile dédiée. Onfleet est reconnue pour
la qualité de son API REST, qui facilite son intégration avec des systèmes tiers, et pour
ses tableaux de bord analytiques permettant de suivre les indicateurs clés de performance
des équipes de livraison.

\subsection{Critique des solutions existantes}

Dans cette section, nous effectuons une critique approfondie des solutions étudiées. Nous
examinons leurs points forts ainsi que leurs limites, afin d'identifier les opportunités
auxquelles notre solution personnalisée doit répondre.

\subsubsection*{Bringgg}

Bringgg propose une plateforme moderne et extensible avec des fonctionnalités avancées
de suivi temps réel. Toutefois, l'intégration avec les ERP utilisés par les entreprises tunisiennes, tels qu'Odoo
ou Dux, nécessite généralement des développements spécifiques ou l'utilisation de solutions
tierces telles que Zapier ou Make, ce qui augmente la complexité et le coût d'intégration.
De plus, Bringgg cible principalement les grandes entreprises avec un modèle tarifaire
sur devis personnalisé, sans grille de prix publique, ce qui le rend peu accessible pour
les PME et les entreprises en croissance.

\subsubsection*{Onfleet}

Onfleet offre une interface ergonomique et de bonnes capacités de suivi opérationnel.
Cependant, l'intégration avec les ERP n'est pas native dans le contexte des solutions
utilisées en Tunisie et repose principalement sur des APIs ou des outils d'automatisation
tiers, ce qui nécessite des développements supplémentaires et une maintenance accrue.
Par ailleurs, son modèle tarifaire basé sur le volume de livraisons peut devenir coûteux
pour les entreprises à forte activité.

Il est important de souligner que la solution que nous souhaitons développer doit pouvoir
se connecter à n'importe quel ERP — qu'il s'agisse d'Odoo, de Dux ou d'un autre système
à l'avenir — sans avoir à tout réécrire à chaque changement. Elle doit également permettre
de gérer plusieurs entreprises clientes depuis une seule plateforme, avec des données
totalement séparées pour chacune.

\subsection{Problématique}

Face aux limites des solutions existantes, les entreprises expriment des besoins clairs
auxquels aucune solution du marché ne répond pleinement :

\begin{itemize}
  \item Savoir en temps réel où se trouvent leurs livraisons et être alertées immédiatement
        en cas de retard ou de dépassement de délai.
  \item Connecter automatiquement leur système de gestion (ERP) à la plateforme de
        livraison, quel que soit l'ERP utilisé, sans ressaisie manuelle des commandes.
  \item Donner à leurs chauffeurs une application mobile simple, utilisable même sans
        connexion internet, pour confirmer les livraisons et capturer des preuves.
  \item Disposer d'une solution qui s'adapte à leur structure, qu'ils aient une seule
        agence ou plusieurs, sans multiplier les abonnements ni les installations.
\end{itemize}

% ─────────────────────────────────────────────────────────────
\section{Solution proposée}

Pour répondre à cette problématique, nous proposons \textbf{ASM Track}, une plateforme SaaS
multi-tenant de gestion des livraisons articulée autour de cinq axes principaux :

\begin{itemize}
  \item \textbf{Architecture microservices} : six services indépendants déployés via Docker
        Compose, communiquant par REST et WebSocket.
  \item \textbf{Intégration ERP flexible} : une architecture permettant de se connecter
        à n'importe quel ERP (Odoo, Dux, ou tout futur système) sans modifier le cœur
        de la plateforme.
  \item \textbf{Fiabilité} : un mécanisme garantissant qu'aucune synchronisation avec
        l'ERP n'est perdue, même en cas de panne réseau ou de redémarrage du serveur.
  \item \textbf{Constructeur de tournées} : un éditeur graphique permettant au dispatcher
        de planifier les tournées par glisser-déposer, avec calcul automatique des
        itinéraires optimaux.
  \item \textbf{Suivi en temps réel} : un tableau de bord pour les dispatchers affichant
        l'état de chaque livraison en direct, avec alertes automatiques en cas de
        dépassement de délai.
  \item \textbf{Application mobile chauffeur} : une application dédiée aux chauffeurs,
        fonctionnelle même sans connexion internet, permettant de confirmer les livraisons,
        capturer des preuves et transférer une tournée à un autre chauffeur via QR code.
  \item \textbf{SaaS multi-tenant} : plusieurs entreprises clientes utilisent la même
        plateforme avec des données totalement séparées, et ASM dispose d'un panneau
        d'administration centralisé pour gérer l'ensemble des clients.
\end{itemize}

% ─────────────────────────────────────────────────────────────
\section{Méthodologie de travail adoptée}

\subsection{Méthodologie Scrum}

Face à la diversité des fonctionnalités à développer et à l'évolution continue des besoins
au cours du projet, il était indispensable de choisir une méthodologie capable de s'adapter
rapidement et de livrer de la valeur de façon régulière. Notre choix s'est porté sur
\textbf{Scrum}, un cadre agile de gestion de projet reconnu pour sa capacité à structurer
le développement logiciel en cycles courts et itératifs, tout en maintenant une vision
claire des priorités.

\subsubsection{Définition de Scrum}

Scrum est un cadre de travail empirique fondé sur trois piliers : la transparence, l'inspection
et l'adaptation. Contrairement aux approches traditionnelles en cascade où tout est planifié
à l'avance, Scrum découpe le développement en cycles courts appelés sprints, à l'issue
desquels un incrément fonctionnel est livré. Pour ASM Track, ce mode de fonctionnement
nous a permis de valider chaque brique technique — authentification, ERP, mobile — avant
de passer à la suivante, sans risquer de construire sur des bases instables.

Le cadre Scrum repose sur trois rôles complémentaires :

\begin{itemize}
  \item Le \textbf{Product Owner} est responsable de définir et de prioriser les
        fonctionnalités du produit en fonction de leur valeur métier. Il maintient
        le Product Backlog à jour et s'assure que l'équipe travaille sur les éléments
        les plus importants à chaque sprint.

  \item Le \textbf{Scrum Master} veille au respect du cadre Scrum et à la bonne
        application de ses pratiques. Il identifie et lève les obstacles qui freinent
        la progression de l'équipe, sans jouer le rôle d'un chef de projet traditionnel.

  \item L'\textbf{équipe de développement} est responsable de la conception, de
        l'implémentation et des tests des fonctionnalités sélectionnées à chaque sprint.
        Elle s'auto-organise pour atteindre l'objectif fixé.
\end{itemize}

Scrum s'appuie également sur cinq valeurs — l'engagement, le courage, le focus,
l'ouverture et le respect — qui ont guidé nos décisions tout au long du projet, notamment
lors des arbitrages techniques entre différentes approches d'implémentation.

\begin{figure}[H]
  \centering
  \includegraphics[width=8cm]{scrum}
  \caption{Les valeurs Scrum}
  \label{fig:scrum-values}
\end{figure}

\subsubsection{Événements Scrum}

Scrum définit cinq événements récurrents qui structurent le travail à l'intérieur de chaque
sprint :

\begin{itemize}
  \item \textbf{Le Sprint} est la période de travail pendant laquelle l'équipe s'engage
        à produire un incrément fonctionnel et potentiellement livrable. Sa durée est
        fixe et ne dépasse pas un mois. Dès la clôture d'un sprint, le suivant commence
        immédiatement, assurant ainsi un rythme de développement continu.

  \item \textbf{Le Sprint Planning} ouvre chaque sprint. L'équipe y sélectionne les
        éléments du Product Backlog à réaliser, estime leur complexité et définit un
        objectif de sprint clair et partagé par tous les membres.

  \item \textbf{Le Daily Scrum} est une réunion quotidienne de quinze minutes destinée
        à synchroniser l'équipe, identifier les obstacles éventuels et ajuster le plan
        de travail de la journée en fonction de l'avancement réel.

  \item \textbf{La Sprint Review} se tient à la fin de chaque sprint. L'incrément
        développé y est présenté aux parties prenantes, dont les retours permettent
        d'affiner le Product Backlog et d'orienter les priorités des sprints suivants.

  \item \textbf{La Sprint Retrospective} clôture chaque sprint sur le plan
        organisationnel. L'équipe y analyse son fonctionnement, identifie les points
        d'amélioration et définit des actions concrètes pour progresser dans le sprint
        suivant.
\end{itemize}

\subsection{Backlog produit}

Le Product Backlog regroupe l'ensemble des fonctionnalités identifiées lors de la phase
d'analyse. Chaque entrée est décrite par son titre, sa description, ses critères d'acceptation,
sa priorité et son estimation en jours :

\begin{longtable}{|C{0.6cm}|L{3cm}|L{3.8cm}|L{4.5cm}|C{1.6cm}|C{1.2cm}|}
\caption{Product Backlog — ASM Track} \\
\hline
\rowcolor{primary!10}
\textbf{N°} & \textbf{Titre} & \textbf{Description} & \textbf{Critères d'acceptation} & \textbf{Priorité} & \textbf{Est. (J)} \\
\hline
\endfirsthead
\hline
\rowcolor{primary!10}
\textbf{N°} & \textbf{Titre} & \textbf{Description} & \textbf{Critères d'acceptation} & \textbf{Priorité} & \textbf{Est. (J)} \\
\hline
\endhead

1 & Serveur OAuth2 &
  Mise en place d'un serveur d'authentification OAuth2 avec JWT signé RSA &
  Les tokens sont émis via /oauth2/token, signés RSA, et validés par l'API Gateway via JWKS &
  Haute & 3 \\
\hline
2 & Connexion admin &
  L'admin se connecte avec son email et mot de passe &
  Token d'accès JWT émis, rôle et identifiant entreprise embarqués &
  Haute & 3 \\
\hline
3 & Connexion chauffeur &
  Le chauffeur se connecte avec son numéro de téléphone &
  Token JWT émis, accès à l'espace mobile autorisé &
  Haute & 3 \\
\hline
4 & Cycle de vie livraison &
  Gestion complète du cycle d'une livraison de la création à la clôture &
  9 états gérés, de UNSCHEDULED jusqu'à DELIVERED, FAILED ou CANCELLED, historique enregistré &
  Haute & 3 \\
\hline
5 & Assignation chauffeur &
  L'admin assigne une livraison à un chauffeur disponible &
  Statut passe à SCHEDULED, notification push FCM envoyée au chauffeur &
  Haute & 3 \\
\hline

6 & Import commandes &
  Les commandes sont importées depuis Odoo ou Dux &
  Les livraisons apparaissent en statut UNSCHEDULED dans le tableau de bord &
  Haute & 4 \\
\hline
7 & Synchronisation ERP &
  Les événements de livraison sont synchronisés vers l'ERP de manière fiable &
  L'ERP est mis à jour automatiquement, idempotence garantie en cas de retry &
  Haute & 4 \\
\hline
8 & Adaptateur multi-ERP &
  Basculer entre Odoo et Dux sans modifier le code métier &
  Le changement d'ERP se fait par configuration, sans redéploiement &
  Haute & 4 \\
\hline

9 & Constructeur de tournées &
  Le dispatcher planifie les tournées par glisser-déposer &
  L'itinéraire est calculé via OSRM (carte Tunisie), les arrêts sont ordonnés &
  Haute & 5 \\
\hline
10 & Suivi temps réel &
  Le dispatcher voit l'état des livraisons et la position des chauffeurs en direct &
  Les mises à jour sont poussées via WebSocket STOMP/RabbitMQ sans rechargement &
  Haute & 5 \\
\hline
11 & Monitoring SLA &
  Détection automatique des livraisons en dépassement de délai &
  Alerte envoyée au dispatcher en moins de 60 secondes &
  Moyenne & 3 \\
\hline

12 & Mode hors-ligne &
  Le chauffeur continue à enregistrer des livraisons sans connexion internet &
  Les actions sont stockées localement (Hive) et synchronisées au retour du réseau &
  Haute & 4 \\
\hline
13 & Preuve de livraison &
  Le chauffeur capture une photo et une signature à la livraison &
  Le fichier est stocké sur MinIO et visible par l'admin &
  Haute & 4 \\
\hline
14 & Transfert QR &
  Le chauffeur transfère une tournée à un collègue via QR code &
  Le chauffeur cible reçoit les stops instantanément &
  Moyenne & 2 \\
\hline

15 & Rapport PDF &
  L'admin génère un rapport PDF de tournée avec les KPIs &
  Le PDF contient 9 indicateurs de performance et est téléchargeable &
  Moyenne & 3 \\
\hline
16 & Multi-tenant SaaS &
  Plusieurs entreprises utilisent la plateforme avec données isolées &
  Chaque entreprise voit uniquement ses données via filtre Hibernate &
  Haute & 4 \\
\hline
17 & Déploiement Docker &
  L'ensemble de la stack est déployé via Docker Compose &
  Tous les services démarrent avec healthchecks et réseau isolé &
  Haute & 4 \\
\hline

\end{longtable}

\subsection{Sprint planning}

Sur la base du Product Backlog établi, nous avons découpé le projet en cinq sprints
successifs, chacun centré sur un axe fonctionnel précis. Ce découpage nous a permis
de valider progressivement les choix techniques et de livrer des incréments testables
à chaque étape :

\begin{table}[H]
\centering
\caption{Réunion de planification des sprints}
\begin{tabularx}{\textwidth}{|C{1.5cm}|X|C{2.4cm}|C{2.4cm}|}
\hline
\rowcolor{primary!10}
\textbf{Sprint} & \textbf{Description} & \textbf{Date début} & \textbf{Date fin} \\
\hline
Sprint 1 & Authentification et cycle de vie des livraisons & 28/02/2025 & 21/03/2025 \\
\hline
Sprint 2 & Intégration ERP agnostique (Odoo / Dux) & 22/03/2025 & 11/04/2025 \\
\hline
Sprint 3 & Constructeur de tournées et suivi temps réel & 12/04/2025 & 02/05/2025 \\
\hline
Sprint 4 & Application mobile chauffeur & 03/05/2025 & 23/05/2025 \\
\hline
Sprint 5 & Reporting, multi-tenancy et déploiement & 24/05/2025 & 30/06/2025 \\
\hline
\end{tabularx}
\end{table}

\subsection{Outil de planification du travail}

Le versioning du code et le suivi des tâches ont été assurés via \textbf{GitLab}.
Au-delà du contrôle de version, GitLab nous a offert un environnement de gestion
de projet complet grâce à plusieurs fonctionnalités intégrées :

\textbf{Suivi intégré des tâches :} GitLab permet de suivre les tâches via des issues
(tickets). Chaque issue peut être associée à des étiquettes, des assignés, des échéances et
des commentaires, offrant ainsi une vue claire de l'état d'avancement du projet.

\textbf{Relation entre les commits et les merge requests :} GitLab permet de lier
directement les commits et les merge requests aux issues. Chaque changement de code est
ainsi traçable jusqu'à une tâche spécifique, ce qui améliore la transparence et la
responsabilité.

\textbf{Tableaux Kanban :} Les tableaux de projet de GitLab permettent de visualiser
le flux de travail sous forme de tableaux Kanban. Nous avons configuré notre tableau avec
les colonnes suivantes : \textit{To do}, \textit{In progress}, \textit{In review} et
\textit{Done}. Cela aide à identifier les goulots d'étranglement et à améliorer le flux
de travail.

\begin{figure}[H]
  \centering
  \includegraphics[width=\textwidth]{issues}
  \caption{Tableau Kanban GitLab — projet ASM Track}
  \label{fig:gitlab-board}
\end{figure}

En résumé, GitLab offre une plateforme unifiée pour la gestion des tâches, le suivi des
commits et les revues de code, rendant le processus de développement plus transparent
et collaboratif.

% ─────────────────────────────────────────────────────────────
\section{Conclusion}

Ce chapitre a constitué le point de départ de notre projet. Nous avons commencé par
présenter l'organisme d'accueil ASM, puis nous avons mené une étude des solutions
existantes sur le marché — Bringgg et Onfleet — dont l'analyse critique nous a permis
d'identifier les limites auxquelles aucune d'elles ne répond pleinement. Sur cette base,
nous avons formulé notre problématique et décrit la solution proposée : ASM Track,
une plateforme SaaS multi-tenant de gestion des livraisons. Enfin, nous avons présenté
la méthodologie Scrum adoptée, avec son backlog produit de dix-sept user stories réparties
sur cinq sprints et gérées via GitLab. Le chapitre suivant présente l'analyse et la
spécification détaillée des besoins fonctionnels et non fonctionnels de la plateforme.
